\section{Числа}

\begin{description}
    \item[] $\N$ --- множество натуральных чисел $(1, 2, 3, \dots)$.
        На $\N$ определены операции сложения и умножения.
    \item[] $\N_0 \coloneqq \{0\} \cup \N$.
    \item[] $\Z$ --- множество целых чисел $(0, \pm 1, \pm 2, \pm 3, \dots)$.
        На $\Z$ добавляется операция вычитания.
    \item[] $\Q$ --- множество рациональных чисел.
        Элементами являются упорядоченные пары $(p, q)$, где $p \in \Z$ и $q \in \N$ --- взаимно простые числа.
        Обозначается как $p/q$. На $\Q$ добавляется операция деления.
    \item[] $\R$ --- множество действительных (вещественных) чисел.
\end{description}

Множество $\Q$ упорядочено, т.е. для любых $x, y \in \Q$ установлено отношение $x \le y$.

Свойства упорядоченности:
\begin{description}
    \item[Рефлексивность:] $x \le x$.
    \item[Антисимметричность:] если $x \le y$ и $y \le x$, то $x = y$.
    \item[Транзитивность:] если $x \le y$ и $y \le z$, то $x \le z$.
\end{description}

Если заведомо $x \neq y$, то определены отношения строгого порядка $x < y$ и $x > y$,
которые обладают только свойством транзитивности.

\begin{theorem}{Принцип Архимеда}{}
    Для любого $c \in \Q$ существует такое $n \in \N$, что $n > c$.
\end{theorem}
